\documentclass[12pt, a4paper]{article}

% Packages
\usepackage[utf8]{inputenc}
\usepackage[T1]{fontenc}
\usepackage[margin=2.5cm]{geometry}
\usepackage{booktabs}
\usepackage{graphicx}
\usepackage{amsmath}
\usepackage{amssymb}
\usepackage{hyperref}
\usepackage{setspace}
\usepackage{parskip}
\usepackage{float}
\usepackage{caption}
\usepackage{subcaption}
\usepackage{array}
\usepackage{xcolor}
\usepackage{titlesec}
\usepackage{fancyhdr}
\usepackage{natbib}

% Header and footer
\pagestyle{fancy}
\fancyhf{}
\rhead{Group BO — SAAM Project 2026}
\lhead{Portfolio Allocation with a Carbon Objective}
\rfoot{Page \thepage}

% Title formatting
\titleformat{\section}{\large\bfseries}{}{0em}{}[\titlerule]
\titleformat{\subsection}{\normalsize\bfseries}{}{0em}{}

% Spacing
\onehalfspacing

\begin{document}

%---------------------------------------------------------------
% COVER PAGE
%---------------------------------------------------------------
\begin{titlepage}
    \centering
    \vspace*{2cm}
    
    {\LARGE \textbf{Portfolio Allocation with a Carbon Objective}}
    
    \vspace{0.5cm}
    {\large Sustainability Aware Asset Management}
    
    \vspace{0.3cm}
    {\large HEC Lausanne — Spring 2026}
    
    \vspace{2cm}
    
    {\large \textbf{Group BO}}
    
    \vspace{0.5cm}
    {\large [Insert group member names here]}
    
    \vspace{2cm}
    
    {\large May 29, 2026}
    
    \vfill
    
    \begin{abstract}
    \noindent
    This report implements a climate-aware asset management strategy for a 
    European equity portfolio. We construct a minimum-variance (MV) portfolio 
    and a value-weighted (VW) benchmark over the out-of-sample period 2014--2025, 
    then introduce three carbon-constrained portfolios that progressively reduce 
    the portfolio carbon footprint using Scope 1 and Scope 2 CO$_2$ emissions. 
    We find that the MV portfolio is naturally 6--10 times greener than the VW 
    benchmark, and that halving the carbon footprint of the VW benchmark can be 
    achieved at very low financial cost. The net zero portfolio (10\% annual 
    reduction) coincides with the 50\% reduction portfolio, suggesting that a 
    standard net zero glide path is less demanding than a one-time 50\% reduction 
    constraint for this European universe.
    \end{abstract}
    
\end{titlepage}

%---------------------------------------------------------------
% TABLE OF CONTENTS
%---------------------------------------------------------------
\tableofcontents
\newpage

%---------------------------------------------------------------
% SECTION 1 — INTRODUCTION
%---------------------------------------------------------------
\section{Introduction}

This report presents the implementation of a climate-aware asset management 
strategy for a European equity portfolio. The project is structured in two parts. 
Part I constructs a minimum-variance (MV) portfolio and compares it to a 
value-weighted (VW) benchmark over the out-of-sample period 2014--2025. Part II 
introduces carbon emission constraints, constructing three additional portfolios 
that progressively reduce the carbon footprint of the investment strategy.

The investment universe consists of European firms drawn from a dataset of 2,545 
companies with available carbon data. We use Scope 1 and Scope 2 CO$_2$ emissions 
combined as our measure of carbon output. All portfolios are rebalanced annually, 
with performance evaluated at the monthly frequency over the period January 2014 
to December 2025 (144 monthly observations).

%---------------------------------------------------------------
% SECTION 2 — DATA AND CLEANING
%---------------------------------------------------------------
\section{Data and Cleaning}

\subsection{Data Sources}

The following datasets are used throughout the project:

\begin{itemize}
    \item \textbf{Monthly total return index} (\texttt{DS\_RI\_T\_USD\_M}): 
    used to compute monthly simple returns
    \item \textbf{Annual market capitalization} (\texttt{DS\_MV\_T\_USD\_Y}): 
    used for value-weighted portfolio construction and carbon footprint computation
    \item \textbf{Monthly market capitalization} (\texttt{DS\_MV\_T\_USD\_M}): 
    used for value-weighted benchmark weights
    \item \textbf{Scope 1 CO$_2$ emissions} (\texttt{DS\_CO2\_SCOPE\_1\_Y}): 
    annual, in tonnes
    \item \textbf{Scope 2 CO$_2$ emissions} (\texttt{DS\_CO2\_SCOPE\_2\_Y}): 
    annual, in tonnes
    \item \textbf{Revenues} (\texttt{DS\_REV\_Y}): annual, in thousands USD, 
    converted to millions
    \item \textbf{Static firm information}: ISIN, name, country, region
\end{itemize}

\subsection{Cleaning Procedure}

\textbf{Price data.} The total return index is loaded at monthly frequency. 
Prices below 0.5 are treated as missing, as they would generate extreme or 
infinite returns. Returns are computed as simple percentage changes:
\[
R_{i,t} = \frac{P_{i,t} - P_{i,t-1}}{P_{i,t-1}}
\]

\textbf{Delisting.} When a firm's price series ends before the end of the 
sample --- indicated by a name suffix such as ``DEAD - DELIST'' in Datastream 
--- a return of $-100\%$ is assigned at the first month after the last valid 
price. All subsequent months remain missing.

\textbf{Carbon and revenue data.} Datastream error strings are replaced with 
missing values. Forward-filling is applied across years within each firm's row 
to handle isolated missing values. Firms with no carbon data at all are excluded.

\textbf{Revenue conversion.} Revenues are divided by 1,000 to convert from 
thousands USD to millions USD, consistent with the carbon intensity formula.

%---------------------------------------------------------------
% SECTION 3 — INVESTMENT SET
%---------------------------------------------------------------
\section{Investment Set Construction}

At the end of each year $Y$, the investment set is defined as the set of firms 
satisfying all of the following criteria:

\begin{enumerate}
    \item \textbf{Region.} The firm belongs to the European region (EUR) as 
    defined in the static dataset.
    
    \item \textbf{Minimum observations.} The firm has at least 36 non-missing 
    monthly return observations over the 10-year estimation window ending at 
    year $Y$.
    
    \item \textbf{Stale prices.} The proportion of months with zero return over 
    the estimation window does not exceed 50\%. Firms exceeding this threshold 
    are excluded as illiquid.
    
    \item \textbf{Valid end-of-year price.} The firm must have a non-missing 
    return in the final month of the estimation window (December of year $Y$).
    
    \item \textbf{Carbon data availability.} The firm must have Scope 1 or 
    Scope 2 emissions data available at end of year $Y$.
    
    \item \textbf{Extreme carbon filter.} Firms whose $E/\text{Cap}$ ratio 
    (emissions divided by market capitalization) exceeds three times the 
    value-weighted market carbon footprint for that year are excluded. This 
    prevents outlier firms from dominating the portfolio carbon footprint. 
    The filter is applied using only information available at end of year $Y$, 
    avoiding look-ahead bias.
\end{enumerate}

The investment set is constructed separately for each rebalancing year from 
2013 to 2024. The number of eligible firms ranges from approximately 350 to 
450 per year after all filters are applied.

%---------------------------------------------------------------
% SECTION 4 — PART I
%---------------------------------------------------------------
\section{Part I --- Standard Portfolio Allocation}

\subsection{Minimum Variance Portfolio}

The minimum variance (MV) portfolio is constructed by solving the following 
optimization problem at the end of each year $Y$:

\begin{align}
    \min_{\alpha_Y} \quad & \alpha_Y^\top \Sigma_Y \alpha_Y \\
    \text{subject to} \quad & \sum_i \alpha_{i,Y} = 1 \notag \\
    & 0 \leq \alpha_{i,Y} \leq 0.05 \quad \text{for all } i \notag
\end{align}

The covariance matrix $\Sigma_Y$ is estimated from 120 months (10 years) of 
monthly returns ending at December of year $Y$. Missing returns within the 
estimation window are replaced by the cross-sectional mean return for that 
month. A small regularization term ($10^{-6} \times I$) is added to the 
covariance matrix to ensure numerical stability.

The weight cap of 5\% per firm is imposed to prevent excessive concentration. 
The optimization is solved using the Sequential Least Squares Programming 
(SLSQP) method from \texttt{scipy.optimize}, with analytic gradients provided 
for computational efficiency.

Between rebalancing dates, weights drift with realized returns:
\[
\alpha_{i,t+1} = \frac{\alpha_{i,t} \times (1 + R_{i,t+1})}{1 + R_{p,t+1}}
\]

\subsection{Value-Weighted Benchmark}

The value-weighted (VW) benchmark weights each firm by its relative market 
capitalization at the end of the previous month:
\[
w_{i,t} = \frac{\text{Cap}_{i,t}}{\sum_j \text{Cap}_{j,t}}
\]

\subsection{Performance Results}

Table~\ref{tab:part1} presents the out-of-sample performance statistics for 
both portfolios over 2014--2025.

\begin{table}[H]
    \centering
    \caption{Portfolio Performance: MV vs VW (2014--2025)}
    \label{tab:part1}
    \begin{tabular}{lcc}
        \toprule
        \textbf{Metric} & \textbf{P\textsubscript{mv}} & 
        \textbf{P\textsubscript{vw}} \\
        \midrule
        Annualized Return       & 6.56\%   & 8.58\%  \\
        Annualized Volatility   & 12.94\%  & 15.87\% \\
        Sharpe Ratio            & 0.35     & 0.41    \\
        Min Monthly Return      & $-$12.09\% & $-$15.36\% \\
        Max Monthly Return      & 13.35\%  & 18.22\% \\
        \bottomrule
    \end{tabular}
    \smallskip
    
    \footnotesize{\textit{Note:} Risk-free rate set to 2\% per annum. 
    Sharpe ratio computed as $(\bar{\mu} - r_f) / \sigma \times \sqrt{12}$.}
\end{table}

The MV portfolio achieves meaningfully lower volatility (12.94\% vs 15.87\%), 
consistent with its objective of minimizing variance. The return is also lower 
(6.56\% vs 8.58\%), reflecting the well-known return-risk trade-off. The Sharpe 
ratio of the MV portfolio (0.35) is slightly below that of the VW benchmark 
(0.41), suggesting that in this sample the variance reduction does not fully 
compensate for the lower return. The cumulative return chart (Figure~\ref{fig:cum_part1}) 
confirms that the MV portfolio grows more steadily, with smaller drawdowns, 
while the VW benchmark achieves higher terminal wealth.

\begin{figure}[H]
    \centering
    \includegraphics[width=0.85\textwidth]{../outputs/cumulative_returns_mv_vw.png}
    \caption{Cumulative returns of the MV and VW portfolios (2014--2025). 
    The MV portfolio (blue) exhibits lower volatility and smaller drawdowns 
    relative to the VW benchmark (orange), at the cost of lower terminal wealth.}
    \label{fig:cum_part1}
\end{figure}

%---------------------------------------------------------------
% SECTION 5 — PART II
%---------------------------------------------------------------
\section{Part II --- Portfolio Allocation with Carbon Emission Reduction}

\subsection{Carbon Metrics (Section 3.1)}

\textbf{Carbon intensity.} The carbon intensity of firm $i$ in year $Y$ is:
\[
CI_{i,Y} = \frac{E_{i,Y}}{\text{Rev}_{i,Y}}
\]
where $E_{i,Y} = \text{Scope1}_{i,Y} + \text{Scope2}_{i,Y}$ is total CO$_2$ 
emissions in tonnes and $\text{Rev}_{i,Y}$ is revenue in millions USD.

\textbf{Weighted Average Carbon Intensity.} The WACI of the portfolio is:
\[
\text{WACI}^P_Y = \sum_i \alpha_{i,Y} \times CI_{i,Y}
\]

\textbf{Carbon Footprint.} The carbon footprint measures emissions attributed 
to the investor per million USD invested:
\[
\text{CF}^P_Y = \sum_i \alpha_{i,Y} \times \frac{E_{i,Y}}{\text{Cap}_{i,Y}}
\]

For the VW portfolio this simplifies to:
\[
\text{CF}^{vw}_Y = \frac{\sum_i E_{i,Y}}{\sum_i \text{Cap}_{i,Y}}
\]

\begin{table}[H]
    \centering
    \caption{Average Carbon Metrics (2013--2024)}
    \label{tab:carbon}
    \begin{tabular}{lcc}
        \toprule
        \textbf{Metric} & \textbf{P\textsubscript{mv}} & 
        \textbf{P\textsubscript{vw}} \\
        \midrule
        Avg WACI (tCO$_2$ / M USD revenue)  & 54.7  & 135.1 \\
        Avg CF (tCO$_2$ / M USD invested)   & 22.0  & 153.3 \\
        \bottomrule
    \end{tabular}
\end{table}

The MV portfolio is naturally significantly greener than the VW benchmark --- 
approximately 6 to 10 times lower in carbon footprint across the sample. This 
is because the MV optimizer tends to select stable, low-volatility firms, which 
are typically in less carbon-intensive sectors such as utilities, healthcare, 
and financials.

\begin{figure}[H]
    \centering
    \includegraphics[width=0.85\textwidth]{../outputs/carbon_metrics.png}
    \caption{WACI (left) and Carbon Footprint (right) for the MV and VW 
    portfolios over 2013--2024. Both metrics show a general downward trend, 
    reflecting real-world improvements in corporate emissions reporting.}
    \label{fig:carbon_metrics}
\end{figure}

\textbf{Top 10 firms driving carbon intensity.} Table~\ref{tab:top10} reports 
the ten firms contributing most to the MV portfolio's weighted average carbon 
intensity, averaged over the full sample.

\begin{table}[H]
    \centering
    \caption{Top 10 Firms by Average CI Contribution (MV Portfolio)}
    \label{tab:top10}
    \begin{tabular}{llc}
        \toprule
        \textbf{Firm Name} & \textbf{ISIN} & 
        \textbf{Avg CI Contribution} \\
        \midrule
        Elia Group                          & BE0003822393 & 14.00 \\
        ICTL HTLS GP                        & GB00BHJYC057 & 5.14  \\
        Koninklijke Vopak                   & NL0009432491 & 5.09  \\
        National Grid                       & GB00BDR05C01 & 3.95  \\
        Devro (delisted 14/04/23)           & GB0002670437 & 2.70  \\
        Coca-Cola HBC                       & CH0198251305 & 2.51  \\
        Mowi                                & NO0003054108 & 2.13  \\
        United Utilities Group              & GB00B39J2M42 & 1.91  \\
        Dignity (delisted 25/05/23)         & GB00BRB37M78 & 1.67  \\
        Teva Pharmaceutical Industries      & US8816242098 & 1.56  \\
        \bottomrule
    \end{tabular}
\end{table}

\subsection{MV Portfolio with 50\% Carbon Constraint --- 
$P^{mv}_{oos}(0.5)$ (Section 3.2)}

We construct a carbon-constrained MV portfolio by solving:

\begin{align}
    \min_{\alpha_Y} \quad & \alpha_Y^\top \Sigma_Y \alpha_Y \\
    \text{subject to} \quad 
    & \text{CF}^P_Y \leq 0.5 \times \text{CF}^{P_{mv}}_{Y} \notag \\
    & \sum_i \alpha_{i,Y} = 1 \notag \\
    & 0 \leq \alpha_{i,Y} \leq 0.05 \quad \text{for all } i \notag
\end{align}

\begin{table}[H]
    \centering
    \caption{Active Investor: $P_{mv}$ vs $P_{mv}(0.5)$}
    \label{tab:active}
    \begin{tabular}{lcc}
        \toprule
        \textbf{Metric} & \textbf{P\textsubscript{mv}} & 
        \textbf{P\textsubscript{mv}(0.5)} \\
        \midrule
        Annualized Return       & 6.56\%   & 7.29\%  \\
        Annualized Volatility   & 12.94\%  & 12.61\% \\
        Sharpe Ratio            & 0.35     & 0.42    \\
        Min Monthly Return      & $-$12.09\% & $-$12.23\% \\
        Max Monthly Return      & 13.35\%  & 13.63\% \\
        Avg CF (tCO$_2$ / M USD) & 22.0   & 10.9    \\
        \bottomrule
    \end{tabular}
\end{table}

Interestingly, $P_{mv}(0.5)$ achieves both higher return and lower volatility 
than the unconstrained $P_{mv}$, resulting in a higher Sharpe ratio (0.42 vs 
0.35). This counterintuitive result arises because the carbon constraint 
effectively forces the optimizer away from certain high-emission, 
higher-volatility firms, redirecting weights toward cleaner, more stable 
companies.

\subsection{Tracking Error Minimization with 50\% Carbon Constraint --- 
$P^{vw}_{oos}(0.5)$ (Section 3.3)}

For the otherwise passive investor, we minimize tracking error relative to the 
VW benchmark subject to a 50\% carbon footprint reduction:

\begin{align}
    \min_{\alpha_Y} \quad & (\alpha_Y - \alpha^{vw}_Y)^\top \Sigma_Y 
    (\alpha_Y - \alpha^{vw}_Y) \\
    \text{subject to} \quad 
    & \text{CF}^P_Y \leq 0.5 \times \text{CF}^{P_{vw}}_{Y} \notag \\
    & \sum_i \alpha_{i,Y} = 1 \notag \\
    & 0 \leq \alpha_{i,Y} \leq 0.05 \quad \text{for all } i \notag
\end{align}

\begin{table}[H]
    \centering
    \caption{Passive Investor: $P_{vw}$ vs $P_{vw}(0.5)$}
    \label{tab:passive}
    \begin{tabular}{lcc}
        \toprule
        \textbf{Metric} & \textbf{P\textsubscript{vw}} & 
        \textbf{P\textsubscript{vw}(0.5)} \\
        \midrule
        Annualized Return       & 8.58\%   & 8.26\%  \\
        Annualized Volatility   & 15.87\%  & 15.97\% \\
        Sharpe Ratio            & 0.41     & 0.39    \\
        Min Monthly Return      & $-$15.36\% & $-$15.42\% \\
        Max Monthly Return      & 18.22\%  & 18.22\% \\
        Avg CF (tCO$_2$ / M USD) & 153.3  & 38.6    \\
        \bottomrule
    \end{tabular}
\end{table}

The carbon constraint reduces the Sharpe ratio slightly (0.39 vs 0.41) and 
marginally increases volatility. However, the financial cost is small --- the 
return difference is only 32 basis points per year. The cumulative return paths 
of $P_{vw}$ and $P_{vw}(0.5)$ are nearly identical, confirming that halving 
the carbon footprint of the VW benchmark can be achieved at very low financial 
cost for a passive European investor.

\subsection{Comparison of Active and Passive Strategies (Section 3.4)}

\begin{figure}[H]
    \centering
    \includegraphics[width=0.85\textwidth]{../outputs/carbon_footprint_comparison.png}
    \caption{Carbon footprint (top) and cumulative returns (bottom) for active 
    investor portfolios (left) and passive investor portfolios (right).}
    \label{fig:comparison}
\end{figure}

The carbon constraint has very different effects depending on the investor type. 
For the active investor, the constraint improves financial performance, 
suggesting that the unconstrained MV portfolio was overweighting certain 
carbon-intensive, relatively volatile firms. For the passive investor, the 
constraint imposes a small but measurable cost in terms of Sharpe ratio, 
consistent with the theoretical prediction that constraints cannot improve the 
unconstrained optimum.

A key finding is that the VW benchmark has a carbon footprint approximately 
7 times higher than the MV portfolio on average. This suggests that European 
equity markets remain significantly carbon-intensive at the market-cap-weighted 
level, and that active low-volatility strategies naturally provide a carbon 
co-benefit without explicit sustainability constraints.

\subsection{Net Zero Portfolio --- $P^{vw}_{oos}(NZ)$ (Section 4.1)}

The net zero portfolio implements a 10\% annual reduction in carbon footprint 
from the 2013 VW baseline:

\[
\text{CF}^P_Y \leq (1 - 0.10)^{Y - 2013} \times \text{CF}^{vw}_{2013}
\]

for $Y = 2013, \ldots, 2024$, where $\text{CF}^{vw}_{2013} = 222.24$ tonnes 
CO$_2$ per million USD invested. The optimization problem is identical to 
Section 3.3 (tracking error minimization) with the net zero constraint 
replacing the 50\% constraint.

\subsection{Comparison of Passive Strategies (Section 4.2)}

\begin{table}[H]
    \centering
    \caption{Passive Investor: $P_{vw}$ vs $P_{vw}(0.5)$ vs $P_{vw}(NZ)$}
    \label{tab:passive_all}
    \begin{tabular}{lccc}
        \toprule
        \textbf{Metric} & \textbf{P\textsubscript{vw}} & 
        \textbf{P\textsubscript{vw}(0.5)} & 
        \textbf{P\textsubscript{vw}(NZ)} \\
        \midrule
        Annualized Return       & 8.58\%     & 8.26\%     & 8.26\%     \\
        Annualized Volatility   & 15.87\%    & 15.97\%    & 15.97\%    \\
        Sharpe Ratio            & 0.41       & 0.39       & 0.39       \\
        Min Monthly Return      & $-$15.36\% & $-$15.42\% & $-$15.42\% \\
        Max Monthly Return      & 18.22\%    & 18.22\%    & 18.22\%    \\
        \bottomrule
    \end{tabular}
\end{table}

$P_{vw}(NZ)$ is identical to $P_{vw}(0.5)$ across all years. This occurs 
because the net zero constraint --- a 10\% annual reduction from the 2013 
baseline of 222.24 tonnes --- is never binding given that the investment set 
already achieves lower carbon footprints than required by the net zero path. 
The 50\% constraint (target $\approx$ 45--111 tonnes depending on the year) 
is consistently tighter than the net zero target (which declines from 200 in 
2014 to 64 by 2024). Since $P_{vw}(0.5)$ already achieves carbon footprints 
in the range of 21--72 tonnes --- well below the net zero target --- the two 
portfolios coincide.

This is an important practical finding: for this European investment universe, 
committing to a one-time 50\% carbon reduction already places the portfolio 
on a trajectory more aggressive than a standard 10\% per year net zero glide 
path.

\begin{figure}[H]
    \centering
    \includegraphics[width=0.85\textwidth]{../outputs/passive_portfolios_comparison.png}
    \caption{Carbon footprint (left) and cumulative returns (right) for 
    passive investor portfolios: $P_{vw}$, $P_{vw}(0.5)$, and $P_{vw}(NZ)$. 
    The net zero target path is shown as a dashed green line.}
    \label{fig:passive_all}
\end{figure}

%---------------------------------------------------------------
% SECTION 6 — LIMITATIONS
%---------------------------------------------------------------
\section{Limitations}

\begin{itemize}
    \item \textbf{Data quality and coverage.} Carbon data coverage is limited 
    before 2010 and uneven across firms, particularly for smaller European 
    companies. Forward-filling of missing emissions data may introduce 
    measurement error in carbon intensity calculations, especially when a 
    firm's actual emissions change significantly year to year.

    \item \textbf{Estimation risk.} The covariance matrix is estimated from 
    120 monthly observations for a universe of 350--450 firms. With more 
    parameters than observations, the sample covariance matrix is poorly 
    conditioned. While regularization and the 5\% weight cap mitigate this, 
    estimation error remains a significant source of out-of-sample performance 
    degradation.

    \item \textbf{Carbon constraint feasibility and outlier sensitivity.} The 
    $E/\text{Cap}$ ratio filter (excluding firms above $3\times$ the market 
    carbon footprint) was necessary to stabilize the carbon footprint of the 
    MV portfolio. While economically justified, this filter reduces the 
    investment universe and may introduce a bias toward smaller, cleaner firms. 
    The choice of $3\times$ as the threshold is somewhat arbitrary.

    \item \textbf{$P_{vw}(NZ)$ and $P_{vw}(0.5)$ coincidence.} The net zero 
    portfolio is identical to the 50\% reduction portfolio throughout the 
    sample. This is because the net zero glide path (10\% annual reduction 
    from 222 tonnes) is less demanding than the 50\% reduction constraint for 
    most of the sample period. A steeper net zero path (e.g., 15\% per year) 
    or a lower baseline would be needed to differentiate the two strategies.

    \item \textbf{Parameter sensitivity.} Key implementation choices --- the 
    5\% weight cap, the 36-observation minimum, the 50\% stale price threshold, 
    and the $3\times$ carbon outlier filter --- all affect results. No 
    sensitivity analysis is conducted on these parameters, and different choices 
    could yield meaningfully different portfolio compositions and performance 
    statistics.
\end{itemize}

%---------------------------------------------------------------
% SECTION 7 — LLM DISCLOSURE
%---------------------------------------------------------------
\section{Use of Large Language Models (LLMs)}

This project made use of Claude (Anthropic) as a support tool during 
implementation. The LLM was used for the following purposes: debugging Python 
code and resolving error messages, understanding pandas and scipy API behaviour, 
improving code structure and readability, and polishing the language and 
structure of this report. All methodological choices --- including the 
investment set filters, the carbon outlier exclusion rule, the choice of SLSQP 
optimization, the weight cap, and the interpretation of all results --- were 
made by the group. All code was written, reviewed, and understood by group 
members. The generation of results, tables, and figures was performed entirely 
by the group's own Python notebooks. The group remains fully responsible for 
the correctness of the implementation and the academic integrity of this 
submission.

\end{document}